\documentclass[16pt]{beamer}
\usepackage[utf8]{inputenc}
\usepackage{pgfpages}
\usepackage{chngpage}
\usepackage{calc}
\RequirePackage[absolute, overlay]{textpos}

\title{All your databases belong to us}
\setbeameroption{show notes on second screen=right}
\usetheme{default}

\usepackage[utf8]{inputenc}
\usepackage{amsmath}
\usepackage{amsfonts}
\usepackage{amssymb}
\usepackage{pgf}
\usepackage{color}
\usepackage[frenchb]{babel}
\usepackage{amssymb}
\usepackage{hyperref}

\usefonttheme{default}
\usepackage{DejaVuSans}
%\usepackage[sfdefault]{FiraSans} %% option 'sfdefault' activates Fira Sans as the default text font
\usepackage[T1]{fontenc}
\renewcommand*\oldstylenums[1]{{\firaoldstyle #1}}


\setbeamertemplate{navigation symbols}{} %remove navigation symbols

\author{Cédric Jeanneret (aka \href{https://www.twitter.com/SwissTengu}{@SwissTengu})}
\institute{\href{https://www.ethack.org/}{EthACK.org}}
\date{\today}

\definecolor{linecolor}{HTML}{4d4c4c}

\setbeamercolor{linecolor}{fg=white,bg=linecolor}

\setbeamertemplate{headline} {
	\begin{beamercolorbox}[wd=\paperwidth,dp=8pt,ht=12pt,leftskip=.29cm,rightskip=.29cm]{linecolor}
	\hfill
	\hypersetup{
		colorlinks=true,
		linkcolor=white,
		urlcolor=white,
	}
	\insertinstitute
	\end{beamercolorbox}%
}

\setbeamertemplate{footline}{%
	\begin{beamercolorbox}[wd=\paperwidth,dp=9pt,ht=0.4cm,leftskip=.29cm,rightskip=.3cm]{linecolor}
	\pgfputat{\pgfxy(0.455,-0.315)}{\pgfbox[center,base]{\includegraphics[width=1.5cm]{../common/logo_537.png}}}
	\hfill
	\inserttitle
	\end{beamercolorbox}%
}


\hypersetup{
	colorlinks=true,
	linkcolor=blue,
	urlcolor=blue,
	pdfborderstyle={/S/U/W 1},
	pdfborder=0 0 1,
	linkbordercolor={0 0 0},
	urlbordercolor={0 0 0},
}


\begin{document}

{
\setbeamertemplate{footline}{%
	\begin{beamercolorbox}[wd=\paperwidth,dp=8pt,ht=12pt,leftskip=.29cm,rightskip=.3cm]{linecolor}
	\hfill
	\inserttitle
	\end{beamercolorbox}%
}

% center first slide — not a title, but almost
{
\centering
\begin{frame}

EthACK
\vspace{0.5cm}

The Swiss Privacy Basecamp 
\vspace{0.5cm}

\includegraphics[width=4cm]{../common/logo_537.png}

\end{frame}
}
}

\begin{frame}{EthACK ?}
\begin{itemize}
	\item Éthique
	\item État
	\item ACKnowledgement (reconnaissance)
	\item Hacking (éthique, évidemment)
	\item …
\end{itemize}
\end{frame}

\begin{frame}{Pourquoi ?}
\begin{itemize}
	\item Notre gouvernement ne s'intéresse pas (ou peu) au sujet
	\item Les sociétés privées nous fichent à notre insu
	\item Personne ne sait où sont leurs données, qui les traitent, à quoi elles servent
\end{itemize}
\end{frame}

\begin{frame}
  \titlepage
\end{frame}

\begin{frame}
{Un sysadmin, c'est...}
\begin{itemize}
\item Le mec barbu qu'on ne voit jamais ou presque
\item Le mec barbu qui râle sur les développeurs
\item Le mec barbu sur lequel les développeurs râlent
\end{itemize}
\end{frame}

\begin{frame}
{Mais pas que ;)}
\begin{itemize}
\item Firewall
\item Updates
\item configuration de services
\item ... etc
\end{itemize}
\end{frame}

\begin{frame}
{Pwnage : surface d'attaque}
\begin{itemize}
\item Service mal configuré
\item Accès utilisateurs trop importants/mal gérés
\item Application en "admin party"
\item Faille de sécurité au niveau de l'application
\item Faille de sécurité au niveau système/OS
\end{itemize}
\end{frame}

\begin{frame}
{Réduction du champ}
\begin{itemize}
\item Application Web
\item Serveur Linux
\item Admin pas réveillé
\item Développeur tout aussi peu réveillé
\end{itemize}
\end{frame}


\begin{frame}
{Compromettre un système d'information}
\begin{itemize}
\item Repérer les applications web installées
\item Repérer les versions
\end{itemize}
\note<1>[item]{README}
\note<1>[item]{Changelog}
\note<1>[item]{Autres fichiers courants (install, etc)}
\end{frame}

\begin{frame}
{Faille détectée : exploitation}
\centering
On est dans la mouise...
\end{frame}

\begin{frame}
{Exemple concret : phpMyAdmin}
\begin{itemize}
\item Très souvent présent
\item Bourré de failles
\item Rarement mis à jour et pris en compte
\end{itemize}
\end{frame}

\begin{frame}
{Dernières "features"}
\begin{itemize}
\item CSRF
\item XSS
\item SQL Injection via nom de BDD
\item Directory transversal (teste l'existence d'un fichier)
\item MitM
\item ......
\end{itemize}
\end{frame}

\begin{frame}
{Autre exemple : Elasticsearch}
\begin{itemize}
\item Service configuré pour écouter par défaut sur toutes les interfaces
\item Aucun mécanisme de sécurité
\end{itemize}
\end{frame}

\begin{frame}
{Conséquences : exploitation possible !}
\begin{itemize}
\item Intégration dans un botnet
\item Infection d'autres hosts du réseau
\item Mapping réseau/infra
\item Base pour attaques ultérieures contre l'infra
\end{itemize}
\end{frame}

\begin{frame}
{Conséquences...}
\begin{itemize}
\item Vols de données
\item Perte de clientèle
\item Pertes de temps
\item Blacklisting d'IPs, coupures réseau
\item La fin du monde des Licornes
\end{itemize}
\end{frame}

\begin{frame}
{Comment éviter ?}
\begin{itemize}
\item Communication (devops-bingo-buzz)
\item Réduire la surface
\item Communication (devops-bingo-buzz-bis)
\item Bonnes pratiques
\end{itemize}
\note<1>[item]{Quels droits, quels ports, quels services}
\note<1>[item]{Enlever les références de version, maintenir à jour, etc}
\note<1>[item]{Cycle des mises à jour}
\note<1>[item]{"Château-fort": mauvais (l'ennemi peut aussi être à l'intérieur); Blacklisting: mauvais (préférer whitelisting)}
\end{frame}

\begin{frame}
{Et les attaques type Déni de service ?}
\begin{itemize}
\item "Scalabilité" de l'application
\item Recherche des bottle-necks
\end{itemize}
\vspace{1cm}
\centering
Tests de charge et communication entre départements !
\note<1>[item]{Note: un Déni de service peut ne pas être dû à une attaque...}
\note<1>[item]{Scalabilité: capacité de l'appli à être distribuée sur plusieurs serveurs}
\note<1>[item]{Stateless, etc, etc}
\end{frame}


{
\setbeamertemplate{footline}{%
	\begin{beamercolorbox}[wd=\paperwidth,dp=8pt,ht=12pt,leftskip=.29cm,rightskip=.3cm]{linecolor}
	\hfill
	\inserttitle
	\end{beamercolorbox}%
}
{
\centering
\begin{frame}
{Questions ?}

\href{https://ethack.org/}{https://ethack.org/} \\
\vspace{0.3cm}
\href{https://www.twitter.com/EthACK_org}{@EthACK\_org} on Twitter \\
\vspace{0.3cm}
\href{https://www.facebook.com/ethack.org}{ethack.org} on Facebook

\vspace{0.5cm}

\includegraphics[width=4cm]{../common/logo_537.png}
\end{frame}
}
}

\end{document}